\documentclass[10pt,a4paper]{amsart}
\usepackage[margin=19mm]{geometry}
\usepackage{amsmath,amssymb,mathtools}
\usepackage[scr=rsfs,cal=euler]{mathalfa}
\usepackage{xfrac}
\usepackage[unicode,hidelinks,bookmarks=false]{hyperref}
\newcommand{\ie}{i.e.\ }
\newcommand{\cf}{cf.\ }
\newcommand{\resp}{respectively\ }
\newcommand{\Subsection}{\subsection}
\providecommand{\nobreakdash}{\nobreak\hbox{-}}
\setlength{\emergencystretch}{2em}
\multlinegap0pt
\usepackage{amssymb}
\usepackage{mathtools}
\usepackage{algorithm}\usepackage{algpseudocode}
\usepackage[shortlabels]{enumitem}
\usepackage{bm}
\usepackage{xcolor}
\usepackage[normalem]{ulem}
\newtheorem{lemma}{Lemma}
\newcommand{\bvec}[1]{\mathbf{#1}}
\title{Dynamic Driver Allocation Under Latent Demand Regimes: Indexability of a Partially Observed Markov Decision Process}
\author{Pedro Cesar Lopes Gerum, Jiong Liu, Ellen Bernal Cavalheiro, Luiz Felipe Martins, Matteo Giaretti}
\date{Source: arXiv:2607.06816v1 (CC BY 4.0)}
\begin{document}
\maketitle

\noindent\emph{Source and licence.} Dynamic Driver Allocation Under Latent Demand Regimes: Indexability of a Partially Observed Markov Decision Process. Version arXiv:2607.06816v1 (CC BY 4.0), distributed by arXiv under the Creative Commons Attribution 4.0 International licence, http://creativecommons.org/licenses/by/4.0/, as recorded in the arXiv licence metadata for this exact version and shown on the abstract page https://arxiv.org/abs/2607.06816v1. Full licence notice: \texttt{licenses/arxiv-2607.06816-CC-BY-4.0.txt}.\par\smallskip

\noindent\textit{Editorial scope.} Counted unit: the article's lemma lem:translation (source0.tex lines 941--949). The article's complete original proof of it is delivered with it (source0.tex lines 951--961, one \texttt{proof} environment of 11 lines). No mathematics is written, repaired or re-derived: the statement and the proof are the article's own lines, in the article's own notation, and no retained line is substituted or re-flowed.  No further source-local context is needed: every cross-reference of every retained line resolves inside the two delivered spans.  Nothing was omitted from any retained span, so the package carries no omission record.  No retained line contains a citation, so the wrapper carries no bibliography and no bibliography entry.  No retained line contains a figure environment, an image-inclusion command or drawing code, so no image is delivered and the package declares no asset dependency.  Editorial formatting only, in addition: an AMS \texttt{amsart} preamble that restates the article's own declarations and macros used by the retained lines, namely the environments \texttt{lemma} on separate counters, together with the standard packages those lines need.  The retained environments print in this document as Lemma 1 in order of appearance, whereas the article prints the same items with the numbering of its own full document.  The complete native source of the article as distributed in the arXiv e-print for this exact version is bundled unchanged as \texttt{sources/204664/source0.tex}.
\begin{lemma}
\label{lem:translation}
Fix a period $t$ and a belief $\bvec{b}$, and assume (A3). For every $s \in \mathbb{Z}_{\ge 0}$ and every $a \in [a_{\min}, a_{\max}-1]$, the $Q$-function satisfies
\begin{equation}
Q_t(s+v,\,a+1) \;=\; Q_t(s,a) \;+\; \kappa, \qquad \kappa \;:=\; (q-c_b)\,v - c_a,
\label{eq:translation}
\end{equation}
with $\kappa < 0$ by (A2).
\end{lemma}

\begin{proof}
Compare the two state-action pairs $(s,a)$ and $(s+v,a+1)$. With $e := x+s$, moving to $(s+v,a+1)$ replaces the effective demand $e$ by $e+v$ and the capacity $va$ by $v(a+1)=va+v$, and the served quantity and next-period backlog transform by elementary algebra. For the served quantity, adding the same constant $v$ to both entries of a minimum shifts the minimum by $v$,
\[
\min\{e+v,\,v(a+1)\} \;=\; \min\{e+v,\,va+v\} \;=\; v + \min\{e,\,va\},
\]
so the served amount rises by exactly $v$. For the next-period backlog, expanding $v(a+1)=va+v$ cancels the added backlog against the added capacity,
\[
\bigl(e+v - v(a+1)\bigr)^+ \;=\; (e+v-va-v)^+ \;=\; (e - va)^+,
\]
so the carried-over backlog, and with it the continuation value, is unchanged. Moving from $(s,a)$ to $(s+v,a+1)$ therefore raises the served-revenue term $q\min\{e,va\}$ by $qv$, raises the backlog penalty $-c_b s$ by $-c_b v$, and raises the wage $-c_a a$ by $-c_a$, while the continuation value is unchanged because $\bvec{b}_{t+1}(x)$ is action-independent under (A3). Summing these increments gives \eqref{eq:translation} with $\kappa = (q-c_b)v - c_a$, and $\kappa < 0$ by (A2).
\end{proof}

\end{document}
